
\lussec 4. 手征对称性关系

由于夸克流方程 (2.4) 保持手征对称性，
像密度 (2.8) 这样的随时间演化复合场
在手征转动下以简单的方式变换，
从而使我们能以有用的新方式研究该对称性的自发破缺。
本节的主要目标，是把随时间演化的凝聚 (2.17)
与手征极限下的普通手征凝聚
及赝标衰变常数联系起来。


\lussubsec 4.1 夸克场方程与 PCAC 关系

现在把时空维数取为 4，
并暂时忽略正规化与重正化的需要，
也就是说，严格地讲本小节的论述只在微扰论的树图级成立。

或许有些出乎意料，
$4+1$ 维理论中的夸克场方程
\equation{
  \langle\{(\slash{D}+M_0)\psi(x)-\lambda(0,x)\}\kern1pt
  \phi_1(t_1,x_1)\ldots\phi_n(t_n,x_n)
  \rangle=
  \hbox{contact terms},
  \enum
}
与 QCD 中的场方程相差一项，
该项与流时间非零处的夸克场耦合。
式 (4.1) 对任意局域场 $\phi_1,\ldots,\phi_n$ 成立，
而当 $4+1$ 维中所有点 $(t_1,x_1),\ldots,(t_n,x_n)$
都不同于 $(0,x)$ 时接触项为零。
场方程 (4.1) 的正确性由恒等式
\equation{
  (\slash{D}+M_0)\wick{\psi(x)\psibar(y)}
  =\longwick{\lambda(0,x)\psibar(y)}+\delta(x-y),
  \enum
  \nexteq{2.5ex}
  (\slash{D}+M_0)\wick{\psi(x)\chibar(s,y)}
  =\longwick{\lambda(0,x)\chibar(s,y)},
  \enum
  \nexteq{2.5ex}
  (\slash{D}+M_0)\wick{\psi(x)\lambdabar(s,y)}
  =\longwick{\lambda(0,x)\lambdabar(s,y)},
  \enum
}
得出，它们是 3.3 小节所引费米子场 Wick 收缩公式的直接推论。
此外，式 (4.2)--(4.4) 表明接触项只来自基本夸克场的收缩。

现在令
\equation{
  \Aax^{rs}_{\mu}(x)=\psibar_r(x)\dirac{\mu}\dirac{5}\psi_s(x),
  \qquad
  \Pax^{rs}(x)=\psibar_r(x)\dirac{5}\psi_s(x),
  \enum
}
为 QCD 的轴矢流与密度。
对于 $r\neq s$，夸克场方程随即蕴含 PCAC 关系
\equation{
  \langle
  \{\partial_{\mu}\Aax^{rs}_{\mu}(x)-(m_{0,r}+m_{0,s})\Pax^{rs}(x)
  +\Ptax^{rs}(x)\}\kern1pt
  \phi_1(t_1,x_1)\ldots\phi_n(t_n,x_n)\rangle
  \noenum
  \nexteq{2.0ex}
  \quad=\hbox{contact terms},
  \enum
}
其中 $m_{0,r},m_{0,s}$ 是味道标记为 $r,s$ 的夸克的裸质量，而
\equation{
  \Ptax^{rs}(x)=\lambdabar_r(0,x)\dirac{5}\psi_s(x)
  +\psibar_r(x)\dirac{5}\lambda_s(0,x).
  \enum
}
特别地，恒等式
\equation{
  \langle
  \{\partial_{\mu}\Aax^{rs}_{\mu}(x)-(m_{0,r}+m_{0,s})\Pax^{rs}(x)
  +\Ptax^{rs}(x)\}\kern1pt P^{sr}_t(y)\rangle=0
  \enum
}
对所有流时间 $t>0$ 以及包括 $x=y$ 在内的所有点 $x$ 成立。
此情形中没有接触项，因为对式 (4.8) 左端关联函数有贡献的
费米子传播子在 $4+1$ 维中全部从 $(0,x)$ 走到 $(t,y)$，
或从 $(t,y)$ 走到 $(0,x)$。
因而收缩 (4.2) 不出现。
鉴于流方程的光滑化性质，
事实上该关联函数的所有部分都是 $x$ 的正规函数。

Wick 收缩结构的另一个纯代数推论是关系式
\equation{
  \int\rmd^4x\,\langle\Ptax^{rs}(x)P^{sr}_t(y)\rangle
  =\langle S^{rr}_t(y)\rangle+\langle S^{ss}_t(y)\rangle,
  \enum
}
将它与 PCAC 关系 (4.8) 相结合，
便得到一个恒等式
\equation{
  (m_{0,r}+m_{0,s})\int\rmd^4x\,\langle\Pax^{rs}(x)P^{sr}_t(y)\rangle
  =-\Sigt^{rr}-\Sigt^{ss}
  \enum
}
它使人们能把随时间演化的夸克凝聚与赝标介子的物理联系起来
（见 4.3 小节）。在推导此式的过程中，
假定夸克质量与流时间均为正，
否则在大 $x$ 处没有边界项、在 $x=y$ 处没有不可积奇异性
就无法得到保证。


\lussubsec 4.2 重正化的 PCAC 关系

超出微扰论树图级后，
式 (4.8)--(4.10) 中的场和夸克质量需要重正化。
这些等式于是成为重正化关联函数之间的关系，
其确切形式依赖于所选的归一化约定。
这里有一些自然的选择，
以下的目的就是把它们明确下来。
讨论隐含地假设了理论的合理正规化，
例如维数正规化或保持大多数对称性的 Wilson 格点表述
[\cite{Wilson}]。与前一样，只考虑味非对角道。

非单态轴矢流与密度作乘法重正化：
\equation{
  \rens{\Aax}{\mu}^{rs}=Z_A\Aax_{\mu}^{rs},
  \qquad
  \ren{\Pax}^{rs}=Z_P\Pax^{rs},
  \enum
}
随时间演化的密度也是如此：
\equation{
  \rens{S}{t}^{rs}=Z_{\chi}S_t^{rs},
  \qquad
  \rens{P}{t}^{rs}=Z_{\chi}P_t^{rs}.
  \enum
}
如前所述，由于正流时间处不存在短距离奇异性，
这些场（包括味对角的那些）按其夸克内容重正化，
即如同 $4+1$ 维中非零距离处
随时间演化的夸克场与反夸克场的乘积那样。

对于场 $\Ptax^{rs}$，
其乘法可重正化性
\equation{
  \ren{\Ptax}^{rs}=\tilde{Z}_P\Ptax^{rs},
  \enum
}
不那么明显，
因为该场量纲为 4，可能与量纲 3 和 4 的其他场混合。
然而涉及拉格朗日乘子场 $\lambda,\lambdabar$
的 Wick 收缩具有这样的性质：
除接触项外，
$\Ptax^{rs}$ 与任何由流时间为零的规范场和夸克场构成的局域场的
两点关联函数为零。
因此被这类场所致的发散加性重正化被排除。
剩下只可能涉及拉格朗日乘子场的场，
但 $\Ptax^{rs}$ 是其中唯一具有量纲 4 及所需对称性质的场，
故加性重正化同样不可能，
于是该场必然作乘法重正化。

现在可以这样选取重正化场与重正化夸克质量 $\mR{r}$
的归一化，使得 PCAC 关系取如下形式：
\equation{
  \langle
  \{\partial_{\mu}\rens{\Aax}{\mu}^{rs}(x)-
  (\mR{r}+\mR{s})\ren{\Pax}^{rs}(x)
  +\ren{\Ptax}^{rs}(x)\}\kern1pt\rens{P}{t}^{sr}(y)\rangle=0.
  \enum
}
注意这一约定只固定了轴矢流 $\rens{\Aax}{\mu}^{rs}$、
乘积 $(\mR{r}+\mR{s})\ren{\Pax}^{rs}$
与密度 $\ren{\Ptax}^{rs}$ 三者之间的相对归一化。
不过后者的归一化可以通过要求方程
\equation{
  \int\rmd^4x\,\langle\ren{\Ptax}^{rs}(x)\rens{P}{t}^{sr}(y)\rangle
  =\langle\rens{S}{t}^{rr}(y)\rangle+\langle\rens{S}{t}^{ss}(y)\rangle
  \enum
}
以及恒等式
\equation{
  (\mR{r}+\mR{s})\int\rmd^4x\,
  \langle\ren{\Pax}^{rs}(x)\rens{P}{t}^{sr}(y)\rangle
  =-\rens{\Sigma}{t}^{rr}-\rens{\Sigma}{t}^{ss}
  \enum
}
（其中 $\rens{\Sigma}{t}^{rr}=Z_{\chi}\Sigma_t^{rr}$）成立而确定。
此后仍需对重正化夸克质量与随时间演化密度施加进一步的归一化条件，
但对质量可以采用某个标准约定，
而随时间演化场的归一化在我们关心的等式中往往相消，
故保持未定。

由式 (4.14)、(4.15) 隐含确定的轴矢流的归一化，
与通常选择的、使轴矢荷取整数值的归一化一致。
有不同办法可以说明这一点，
也许最直接的途径是从保持手征对称性的格点 QCD 表述出发
[\cite{GinspargWilson}--\cite{ExactChSy}]。
利用手征 Ward 恒等式 [\cite{BoulderReview}]，
可以证明在此情形下裸场满足式 (4.8)、(4.9)，
直至正比于格距而消失的项。
于是式 (4.14)、(4.15) 意味着 $Z_A=\tilde{Z}_P=1$，
从而给出重正化轴矢流的 canonical 归一化。
鉴于连续极限的普适性，
这一 canonical 归一化因此由这些等式保证，
而与所选的理论正规化无关。


\lussubsec 4.3 赝标矩阵元与手征凝聚

现在假定夸克多重态包含两个轻夸克，
称为 $u$、$d$ 夸克，二者具有相同的
重正化质量 $\ren{m}$。
若再无太多其余夸克，
$(u,d)$ 道的手征对称性预期会发生自发破缺，
于是存在一个赝标介子三重态——``$\pi$ 介子''，
其质量随 $\ren{m}\to0$ 而趋于零。

在大距离处，轻夸克赝标关联函数
由中间单 $\pi$ 介子态主导。
特别地，当 $x_0\to\infty$ 时
\equation{
  \int\rmd^3{\mib x}\,\langle\ren{\Pax}^{ud}(x)\ren{\Pax}^{du}(0)\rangle
  =-{\Gpi^2\over\Mpi}
  \rme^{-\Mpi x_0}\bigl\{1+\rmO(\rme^{-\Delta Ex_0})\bigr\},
  \enum
  \nexteq{2.5ex}
  \int\rmd^3{\mib x}\,\langle\rens{\Aax}{0}^{ud}(x)\ren{\Pax}^{du}(0)\rangle
  =\Fpi\Gpi\rme^{-\Mpi x_0}\bigl\{1+\rmO(\rme^{-\Delta Ex_0})\bigr\},
  \enum
}
其中 $\Delta E$ 表示到更高态的能量间隙，
而 $\Mpi$、$\Fpi$ 与 $\Gpi$
分别是 $\pi$ 介子质量、衰变常数
与轴矢密度的真空到 $\pi$ 介子矩阵元。

从四维理论的观点看，
$\rens{\Pax}{t}^{du}(x)$ 不是局域场，
但由于光滑核 $K(t,x;0,y)$ 在大距离 $|x-y|$ 处
近似按高斯型衰减，
它仍然是基本场的一个相当局域的表达式。
当 $x_0$ 远大于光滑范围 $\sqrt{8t}$ 时，
两点函数
\equation{
  \int\rmd^3{\mib x}\,\langle\ren{\Pax}^{ud}(x)\rens{\Pax}{t}^{du}(0)\rangle
  =-{\Gpi\Gpit\over\Mpi}
  \rme^{-\Mpi x_0}\bigl\{1+\rmO(\rme^{-\Delta Ex_0})\bigr\}
  \enum
}
因而与普通赝标关联函数按同样方式衰减。
系数 $\Gpit$ 等于某个算符的真空到 $\pi$ 介子矩阵元，
但此处把它视为由式 (4.19) 定义的量。

表征（两味）手征极限的低能有效常数
\equation{
  \Sigma=\lim_{m_{\rm R}\to0}\Fpi\Gpi,
  \enum
  \nexteq{2.5ex}
  F=\lim_{m_{\rm R}\to0}\Fpi,
  \enum
}
是手征凝聚 $\Sigma$
以及在 $\ren{m}=0$ 处的 $\pi$ 介子衰变常数值 $F$。
如前所述，
通过手征 Ward 恒等式 (4.16) 可以把它们与随时间演化的
轻夸克凝聚 $\rens{\Sigma}{t}=\rens{\Sigma}{t}^{uu}$
联系起来。利用 PCAC 关系
\equation{
  2\ren{m}\Gpi=\Mpi^2\Fpi,
  \enum
}
后者可以写成如下形式：
\equation{
  \rens{\Sigma}{t}=-{\Mpi^2\Fpi\over2\Gpi}
  \int\rmd^4x\,\langle \ren{P}^{ud}(x)\rens{P}{t}^{du}(0)\rangle.
  \enum
}
此式右端的积分在 $\ren{m}\to0$ 时发散，
其渐近行为由动量空间中关联函数的 $\pi$ 极点决定。
回忆式 (4.19)，即可得出结论
\equation{
  \Sigma=\lim_{m_{\rm R}\to0}{\rens{\Sigma}{t}\Gpi\over\Gpit},
  \enum
  \nexteq{2.5ex}
  F=\lim_{m_{\rm R}\to0}{\rens{\Sigma}{t}\over\Gpit}.
  \enum
}
于是在手征极限下，
随时间演化的凝聚 $\rens{\Sigma}{t}$
提供了手征对称性自发破缺的一个序参量，
它与标准凝聚 $\Sigma$ 只差一个比例常数。

式 (4.24)、(4.25) 的一个显著特点是：
其中完全没有提及轴矢流，
甚至没有通过归一化条件隐含涉及；
而且这些等式对任何（正的）流时间 $t$ 都成立。
此外，重正化常数 $Z_{\chi}$ 消失了，
因此例如衰变常数 $F$
可以直接由裸的随时间演化凝聚
与赝标真空到 $\pi$ 介子矩阵元给出。


\lussubsec 4.4 手征微扰论

在手征极限附近，
QCD 中许多量的渐近行为可以用手征微扰论描述。
如果 $\pi$ 介子的康普顿波长远大于梯度流的光滑范围，即若
\equation{
  8t\Mpi^2\ll1,
  \enum
}
则随时间演化的密度 $\rens{S}{t}^{rs}$、$\rens{P}{t}^{rs}$
从 $\pi$ 介子物理的观点看是局域场，
具有与流时间为零时的密度相同的对称性质。
因此在低能下，
随时间演化场的关联函数以及普通轴矢流与密度的关联函数
预期由 $\pi$ 介子中间态主导，
从而应当可以在手征微扰论框架内得到定量描述。

事实上，只要在手征拉氏量中加入适当的源项
[\cite{ChPT}]，
就可以很容易地把随时间演化的密度纳入手征微扰论。
在手征展开的每一阶，这些项要求引入若干有效耦合，
它们一般会依赖流时间。
于是可以算出手征凝聚 $\rens{\Sigma}{t}$
与矩阵元 $\Gpit$ 的展开，
从而洞察手征极限 (4.24)、(4.25)
究竟是如何达到的。
